\chapter*{Mapa de prerrequisitos y método de estudio}
\addcontentsline{toc}{chapter}{Mapa de prerrequisitos y método de estudio}

\begin{center}
\begin{tikzpicture}[
  >=Latex,
  node distance=9mm and 11mm,
  box/.style={
    rounded corners=2mm,
    minimum width=3.15cm,
    minimum height=1.05cm,
    align=center,
    font=\sffamily\small,
    line width=.75pt
  },
  base/.style={box,draw=BookGray!70,fill=BookGrayLight},
  core/.style={box,draw=BookBlue,fill=BookBlueLight},
  constit/.style={box,draw=BookGreen,fill=BookGreenLight},
  laws/.style={box,draw=BookOrange,fill=BookOrangeLight},
  arr/.style={-{Latex[length=2.4mm]},line width=.9pt,draw=BookBlue!85}
]
\node[base] (stat) {Estática\\fuerzas y momentos};
\node[core,right=of stat] (model) {Modelización\\barra y sección};
\node[core,right=of model] (int) {Esfuerzos internos\\$N,\ V,\ M,\ T$};

\node[constit,below=of model] (stress) {Tensión\\$\sigma,\ \tau$};
\node[constit,left=of stress] (strain) {Deformación\\$\varepsilon,\ \gamma$};
\node[constit,right=of stress] (hooke) {Ley constitutiva\\material elástico};

\node[laws,below=of int] (qvm) {Leyes de esfuerzos\\$V'=-q,\quad M'=V$};

\draw[arr] (stat)--(model);
\draw[arr] (model)--(int);
\draw[arr] (int)--(stress);
\draw[arr] (stress)--(strain);
\draw[arr] (strain)--(hooke);
\draw[arr] (int)--(qvm);
\draw[arr] (stat) |- (qvm);

\node[font=\sffamily\scriptsize,text=BookGray,below=3mm of strain]
  {cinemática};
\node[font=\sffamily\scriptsize,text=BookGray,below=3mm of hooke]
  {comportamiento};
\end{tikzpicture}
\end{center}

\begin{tipbox}
Para cada capítulo sigue cuatro pasadas: \textbf{(1)} lectura física, \textbf{(2)} derivación con lápiz, \textbf{(3)} problema sin mirar, \textbf{(4)} auditoría de signos/unidades. Este patrón evita confundir reconocimiento de una solución con dominio real.
\end{tipbox}

\section*{Método universal de resolución}
El curso puede organizarse alrededor de cuatro principios: diagrama de cuerpo libre, equilibrio, compatibilidad de desplazamientos y relación tensión--deformación. TU Delft presenta explícitamente esta estructura como marco común para los problemas de mecánica de materiales.\autocite{tudelft-four-principles}
